\chapter{CONCLUSION}

This thesis investigated neural network-based axle detection for Bridge
Weigh-In-Motion (B-WIM) systems. The problem was framed as per-sample binary
sequence labelling on a dataset of 32,141 vehicle crossing records, each a
1,300-sample strain signal with ground-truth axle positions. The extreme class
imbalance (287:1 negative-to-positive ratio) was addressed by weighting the
binary cross-entropy loss proportionally to the class frequencies.

Three approaches were compared: a classical Savitzky-Golay filter and
peak-detection baseline, a 1D U-Net CNN (30\,M parameters), and a TCN with nine
dilated residual blocks (2.7\,M parameters, receptive field 2,045 samples). Both
deep learning models were evaluated using axle-level precision, recall, F1-score,
and Mean Absolute Timing Error (MATE) with a $\pm$5-sample tolerance window.

The main findings are:

\begin{enumerate}
    \item Both neural network models achieve near-perfect axle detection
          (F1\,$\geq$\,0.997, MATE\,$\leq$\,0.25 samples) on unseen vehicle
          records, substantially outperforming the classical baseline.
    \item The TCN outperforms the larger CNN on all metrics despite having 11
          times fewer parameters. This demonstrates that \emph{architecture
          design} --- specifically, ensuring sufficient receptive field coverage
          --- matters more than model scale for this task.
    \item A MATE of approximately 0.24 samples corresponds to sub-millisecond
          timing accuracy at 100\,Hz, well within the tolerance required for
          accurate GVW estimation.
\end{enumerate}

\section{FUTURE WORK}

Several directions remain open for future investigation:

\begin{itemize}
    \item \textbf{Day-level data splitting}: Re-evaluating models with a
          chronological split would test cross-temporal generalisation more
          rigorously.
    \item \textbf{Multi-installation generalisation}: Collecting data from
          multiple bridges and sensor configurations would reveal how well the
          models transfer to unseen installations.
    \item \textbf{Robustness under noise and speed variation}: Systematic
          experiments with synthetic noise injection and time-stretching would
          quantify model robustness.
    \item \textbf{End-to-end weight estimation}: Integrating the axle detector
          into a full B-WIM pipeline and evaluating GVW accuracy would demonstrate
          operational impact.
    \item \textbf{Lighter architectures}: Exploring smaller TCN variants or
          knowledge distillation could enable deployment on low-power edge
          hardware typically used in B-WIM installations.
\end{itemize}
